% 推理框架生态架构图 —— 双路径分层架构
% 左侧「轻量路径」(个人/Mac) ←→ 右侧「生产路径」(团队/NVIDIA)
\documentclass[border=8pt,tikz]{standalone}
\usepackage{tikz}
\usetikzlibrary{shapes.geometric, arrows.meta, positioning, fit, backgrounds, calc, decorations.pathreplacing}
\usepackage{xeCJK}
\setCJKmainfont{STSong}

\begin{document}
\begin{tikzpicture}[
    node distance=1.4cm,
    % ---- 颜色体系 ----
    % 轻量路径 (绿/青)
    lw_front/.style={rectangle, rounded corners=4pt, draw=teal!70!black, fill=teal!8,
                     minimum width=2.6cm, minimum height=0.85cm,
                     font=\footnotesize\bfseries, align=center},
    lw_engine/.style={rectangle, rounded corners=4pt, draw=teal!80!black, fill=teal!18,
                      minimum width=2.8cm, minimum height=0.9cm,
                      font=\footnotesize\bfseries, align=center},
    lw_hw/.style={rectangle, rounded corners=3pt, draw=teal!60!black, fill=teal!5,
                  minimum width=2.6cm, minimum height=0.8cm,
                  font=\footnotesize\bfseries, align=center},
    % 生产路径 (蓝/靛)
    prod_front/.style={rectangle, rounded corners=4pt, draw=blue!70!black, fill=blue!8,
                       minimum width=2.8cm, minimum height=0.85cm,
                       font=\footnotesize\bfseries, align=center},
    prod_engine/.style={rectangle, rounded corners=4pt, draw=blue!80!black, fill=blue!22,
                        minimum width=3.0cm, minimum height=0.9cm,
                        font=\footnotesize\bfseries, align=center},
    prod_engine_sec/.style={rectangle, rounded corners=4pt, draw=blue!60!black, fill=blue!8,
                            minimum width=2.8cm, minimum height=0.85cm,
                            font=\footnotesize\bfseries, align=center},
    prod_hw/.style={rectangle, rounded corners=3pt, draw=blue!70!black, fill=blue!5,
                    minimum width=2.8cm, minimum height=0.8cm,
                    font=\footnotesize\bfseries, align=center},
    % 通用
    yaxis_label/.style={font=\small\bfseries, text=black!55, align=center, anchor=center},
    arrow_lw/.style={-{Stealth[scale=0.8]}, line width=0.8pt, draw=teal!50},
    arrow_prod/.style={-{Stealth[scale=0.8]}, line width=0.8pt, draw=blue!50},
    arrow_cross/.style={-{Stealth[scale=0.7]}, line width=0.5pt, draw=gray!45, dashed},
]

% ============================================================
% 硬件后端层 (底层) — 两大列
% ============================================================
% ---- 左侧: 轻量硬件 ----
\node[lw_hw] (apple)  {Apple Silicon\\(Metal / UMA)};
\node[lw_hw] (cpu)    [below=0.5cm of apple]  {x86 CPU\\(AVX2, 回退方案)};

% ---- 右侧: 生产硬件 ----
\node[prod_hw] (nvidia) [right=5.0cm of apple] {NVIDIA GPU\\(CUDA / TensorRT)};

% 左侧 Y 轴标签: 硬件层
\node[yaxis_label, rotate=90]
    at ($(apple.west)!0.5!(cpu.west) + (-1.2,0)$) {硬件后端层};

% ============================================================
% 推理引擎层 (中层)
% ============================================================
% ---- 左侧: 轻量引擎 ----
\node[lw_engine] (ollama)    [above=1.8cm of apple] {Ollama\\\footnotesize 封装 llama.cpp};
\node[lw_engine] (llamacpp)  [above=0.5cm of ollama]    {llama.cpp\\\footnotesize GGUF / CPU-GPU混合};

% ---- 右侧: 生产引擎 ----
\node[prod_engine] (vllm)       [above=1.8cm of nvidia] {vLLM\\\footnotesize PagedAttention};
\node[prod_engine] (sglang)     [left=0.4cm of vllm]         {SGLang\\\footnotesize RadixAttention};
\node[prod_engine_sec] (trtllm) [right=0.4cm of vllm]        {TensorRT-LLM\\\footnotesize 图优化 / FP8};
\node[prod_engine_sec] (tgi)    [right=0.4cm of trtllm]      {HuggingFace TGI\\\footnotesize 连续批处理};

% Ollama→llama.cpp 包含关系
\draw[arrow_lw, thick] (ollama.north) -- (llamacpp.south)
    node[midway, right, font=\footnotesize\itshape, text=teal!60] {封装};

% 左侧 Y 轴标签: 引擎层
\node[yaxis_label, rotate=90]
    at ($(llamacpp.west)!0.5!(ollama.west) + (-0.8,0)$) {推理引擎层};

% ============================================================
% 前端 / UI 层 (顶层)
% ============================================================
% ---- 左侧: 轻量前端 ----
\node[lw_front] (lmstudio)   [above=0.8cm of llamacpp] {LM Studio\\\footnotesize 桌面GUI};
\node[lw_front] (continue)   [above=0.5cm of lmstudio]    {Continue.dev\\\footnotesize IDE 插件};

% ---- 右侧: 生产/团队前端 ----
\node[prod_front] (openwebui) [above=1.8cm of vllm] {Open WebUI\\\footnotesize 团队 Web 界面};
\node[prod_front] (anything)  [left=0.4cm of openwebui]      {AnythingLLM\\\footnotesize 内置 RAG};
\node[prod_front] (ooba)      [right=0.4cm of openwebui]     {text-gen-webui\\\footnotesize 全功能 WebUI};

% 左侧 Y 轴标签: 前端层
\node[yaxis_label, rotate=90]
    at ($(continue.west)!0.5!(lmstudio.west) + (-1.2,0)$) {前端 / Web UI 层};

% ============================================================
% 连接箭头: 前端 → 引擎层
% ============================================================
% 轻量路径 (solid teal)
% \draw[arrow_lw] (continue.south)  -- (llamacpp.north);
% \draw[arrow_lw] (continue.south)  -- (ollama.north);
% \draw[arrow_lw] (lmstudio.south)  -- (ollama.north);

% 生产路径 (solid blue)
\draw[arrow_prod] (anything.south)  -- (vllm.north);
\draw[arrow_prod] (openwebui.south) -- (vllm.north);
\draw[arrow_prod] (openwebui.south) -- (sglang.north);
\draw[arrow_prod] (ooba.south)      -- (tgi.north);
\draw[arrow_prod] (ooba.south)      -- (vllm.north);

% 跨路径连接 (dashed gray) — Open WebUI 也能用 Ollama
\draw[arrow_cross] (openwebui.south west) -- (ollama.north east);

% ============================================================
% 连接箭头: 引擎 → 硬件
% ============================================================
% llama.cpp 跨硬件 (唯一支持三平台的引擎)
% \draw[arrow_lw] (llamacpp.south) -- (apple.north);
% \draw[arrow_lw] (llamacpp.south) -- (cpu.north);
\draw[arrow_cross] (llamacpp.south east) -- (nvidia.north west);

% Ollama → Apple / CPU
% \draw[arrow_lw] (ollama.south) -- (apple.north);
% \draw[arrow_lw] (ollama.south) -- (cpu.north);

% 生产引擎 → NVIDIA (全部仅支持 CUDA)
\draw[arrow_prod] (sglang.south)  -- (nvidia.north);
\draw[arrow_prod] (vllm.south)    -- (nvidia.north);
\draw[arrow_prod] (trtllm.south)  -- (nvidia.north);
\draw[arrow_prod] (tgi.south)     -- (nvidia.north);

% ============================================================
% 推荐路径标注 (绿色虚线框)
% ============================================================
% 推荐路径 1: Continue.dev + Ollama/llama.cpp + Apple Silicon
\node[draw=green!55!black, very thick, dashed, rounded corners=8pt,
      fit=(continue)(lmstudio)(ollama)(llamacpp)(apple)(cpu),
      inner sep=6pt] (rec_personal) {};
\node[green!55!black, font=\footnotesize\bfseries, anchor=north west,
      , fill opacity=0.85, text opacity=1]
    at ($(rec_personal.north east)+(0.3,0.2)$) {$\leftarrow$ 个人 / Mac 推荐};

% 推荐路径 2: Open WebUI / AnythingLLM + vLLM / SGLang + NVIDIA
\node[draw=green!55!black, very thick, dashed, rounded corners=8pt,
      fit=(anything)(openwebui)(sglang)(vllm)(nvidia),
      inner sep=3pt] (rec_prod) {};
\node[green!55!black, font=\footnotesize\bfseries, anchor=north west,
      , fill opacity=0.85, text opacity=1]
    at ($(rec_prod.north east)+(0.3,0.4)$) {$\leftarrow$ 团队 / 生产推荐};

% ============================================================
% 图例
% ============================================================
\node[font=\footnotesize, anchor=north west, align=left, text=black!55]
    at ($(current bounding box.south west)+(0,-0.6)$) {
    \textbf{图例：}
    \tikz{\node[lw_engine, minimum width=0.7cm, minimum height=0.35cm, font=\tiny] {};} 轻量/个人路径 \quad
    \tikz{\node[prod_engine, minimum width=0.7cm, minimum height=0.35cm, font=\tiny] {};} 生产/团队路径 \quad
    --- 主连接 \quad
    -\,-\,- 跨路径兼容 \quad
    \tikz{\draw[green!55!black, very thick, dashed] (0,0) -- (0.8,0);} 推荐组合
};

\end{tikzpicture}
\end{document}
